\documentclass[exercice,a4]{thedocument}%
\usepackage[french]{babel}%
\usepackage{pstricks}%
\usepackage{tabularx}%
\usepackage{graphicx}%
\usepackage{array}%
\startdatakeys
    \source{21GENMATMEAG3}
    \BO{2026}
    \cycle{4}
    \competencesDuSocle{CA,RA}
    \objectifsApprentissage{C4D1be,C4D2aa}
    \motsCles{étendue,moyenne,médiane,pourcentage,coordonnées géographiques}
\enddatakeys
\begin{document}%


\medskip

Sur l'île de Madagascar, un scientifique mène une étude sur les tortues vertes.

La tortue verte a pour nom scientifique: \og Chelonia Mydas \fg.

La carapace mesure en moyenne $115$~cm et l'animal pèse entre $80$ et $130$~kg.

Elle est classée comme espèce \og En Danger».

\medskip

Afin de surveiller la bonne santé des tortues, elles sont régulièrement pesées. Voici les données relevées par ce scientifique en mai 2021.

\begin{center}
\begin{tabularx}{\linewidth}{|m{3cm}|*{7}{>{\centering \arraybackslash}X|}}\hline
Lettres de marquage&A-001&A-002&A-003&A-004&A-005&A-006&A-007\\ \hline
Sexe de la tortue&Mâle&Femelle&Femelle&Femelle&Mâle& Femelle&Femelle\\ \hline
Masse (en kg)&113&96&125&87&117&104&101\\ \hline
\end{tabularx}
\end{center}

\begin{enumerate}
\item Calculer l'étendue de cette série statistique.
\item Calculer la masse moyenne de ces 7 tortues. Arrondir le résultat à l'unité.
\item Déterminer la médiane de cette série statistique. Interpréter le résultat.
\item Est-il vrai que les mâles représentent moins de 20\,\% de cet échantillon ?
\item L'île de Madagascar a pour coordonnées géographiques (20 Sud ; 45 Est).

Placer une croix sur le planisphère fourni en annexe afin de marquer la position de l'île de Madagascar.
\smallskip

L'annexe est à rendre avec la copie.
\end{enumerate}

\medskip

\begin{center}
\vspace{3cm}

\begin{pspicture}(15,10)
%\psgrid
\rput(7.5,5){\eps[width=14.cm]{Mappemonde_Metropole_sept_2021}}
\psdots[linecolor=yellow,dotscale=1.](2.8,6.7)(7.22,7.18)(8.9,4)(12.2,6.675)
\end{pspicture}
\end{center}

\showthesource
\end{document}
